\begingroup\setbox0=\vbox{\section*{Decimals}}\endgroup
\chapterhead{3}{Decimals}

\shead{Decimal fractions -- Decimals}
\begin{knowbox}\begin{itemize}[leftmargin=1.4em,itemsep=3pt]
\item A \textbf{decimal fraction} has denominator $10,\ 100,\ 1\,000,\ \dots$ \ For example $\dfrac4{10},\ \dfrac{45}{100},\ \dfrac{358}{1\,000}$.
\item Every decimal fraction can be written as a \textbf{decimal}: $\ \dfrac4{10}=0.4;\quad\dfrac{45}{10}=4.5;\quad\dfrac{358}{1\,000}=0.358.$
\item Many other fractions can be changed into decimal fractions and then into decimals: $\ \dfrac12=\dfrac5{10}=0.5;\quad\dfrac{13}5=\dfrac{26}{10}=2.6.$
\item A decimal has a \textbf{whole-number part} (left of the decimal point) and a \textbf{decimal part} (right of it). The decimal places are: tenths, hundredths, thousandths,\,\dots
\item One unit of any place equals $10$ units of the next place to the right. Writing (or removing) zeros at the right end of the decimal part does not change the value: $\ 2.5=2.50=2.500$.
\end{itemize}\end{knowbox}
\gap
\shead{Comparing two decimals}
\begin{knowbox}
\begin{itemize}[leftmargin=1.4em,itemsep=3pt]
\item Compare the \textbf{whole-number parts} first: the larger whole part gives the larger number.
\item If the whole parts are equal, compare the decimal places in turn (tenths, hundredths, thousandths,\,\dots). At the first place where the digits differ, the larger digit gives the larger number.
\end{itemize}
Example: $\ 7.305>7.29\ $ because the whole parts are equal and in the tenths $3>2$.
\end{knowbox}
\gap
\shead{Moving the decimal point}
\begin{knowbox}\begin{itemize}[leftmargin=1.4em,itemsep=3pt]
\item Moving the point $1,\ 2,\ 3,\ \dots$ places to the \textbf{right} makes the number $10,\ 100,\ 1\,000,\ \dots$ \textbf{times larger}: $\ 3.75\times10=37.5$.
\item Moving the point $1,\ 2,\ 3,\ \dots$ places to the \textbf{left} makes it $10,\ 100,\ 1\,000,\ \dots$ \textbf{times smaller}: $\ 52.1:100=0.521$.
\item So when the point moves one place left, if the new number is $1$ part, the old number is $10$ such parts: the number has \emph{decreased} by $9$ parts.
\end{itemize}\end{knowbox}
\gap
\begin{exbox}{\ExWord\ 1}
Write these decimals:
\begin{cols3}
\pa{a} $4$ tens, $7$ ones, $5$ tenths, $3$ thousandths; & \pa{b} $5$ ones and $38$ hundredths; & \pa{c} $0$ ones and $7$ thousandths.
\end{cols3}
\guide
Put each digit in its place; write $0$ in any empty place:
\[\text{\pa{a}}\ 47.503;\qquad\text{\pa{b}}\ 5.38;\qquad\text{\pa{c}}\ 0.007.\]
\end{exbox}
\gap
\begin{exbox}{\ExWord\ 2}
Write as fractions in lowest terms: $\ 12.25\ $ and $\ 3.125$.
\guide
Write as a decimal fraction, then simplify:
\[12.25=\frac{1\,225}{100}=\frac{49}{4};\qquad 3.125=\frac{3\,125}{1\,000}=\frac{25}{8}.\]
\end{exbox}
\gap
\begin{exbox}{\ExWord\ 3}
Find all natural numbers $x$ with $\ 2.46<x<\dfrac{31}5$.
\guide
Change $\dfrac{31}5=\dfrac{62}{10}=6.2$.\par We need $2.46<x<6.2$, so $x\in\{3;\ 4;\ 5;\ 6\}$.
\end{exbox}
\gap
\begin{exbox}{\ExWord\ 4}
When the decimal point of a number is moved one place to the left, the number decreases by $27.18$. Find the number.
\guide
Moving the point one place left makes the number $10$ times smaller.\par Bar model:
\begin{barfigC}
\rowlab{-0.15}{0.62}{Old number:}
\foreach \i in {0,...,9}{\blk{\i*0.55}{\i*0.55+0.55}{0.62}{0.42}{tmBlue}{}}
\rowlab{-0.15}{0}{New number:}
\blk{0}{0.55}{0}{0.42}{tmGreen}{}
\brcdn{0.55}{5.5}{-0.32}{$27.18$}
\end{barfigC}
Difference in parts: $10-1=9$.\par One part (the new number): $27.18:9=3.02$.\par
The number: $3.02\times10=30.2$.\par \emph{Check:} $30.2-3.02=27.18$.
\end{exbox}
\gap
\begin{exbox}{\ExWord\ 5}
Calculate quickly: $\ S=0.2+0.4+0.6+\dots+3.8+4$.
\guide
This is an arithmetic sequence with difference $0.2$.\par Number of terms: $(4-0.2):0.2+1=20$.\par
Sum: $S=(0.2+4)\times20:2=42$.
\end{exbox}
\dividerbar{Practice}
\begin{probbox}{1}{\relax}
Write as decimals:
\begin{cols2}
\pa{a} $\dfrac7{10};\ \ \dfrac9{1\,000};\ \ \dfrac{63}{100};\ \ \dfrac{1\,234}{100};\ \ 4\dfrac3{10};\ \ 12\dfrac{45}{1\,000}$ & \pa{b} $\dfrac25;\ \ \dfrac74;\ \ 3\dfrac9{20};\ \ \dfrac{11}8$
\end{cols2}
\end{probbox}
\probsol{
\pa{a} $0.7;\ \ 0.009;\ \ 0.63;\ \ 12.34;\ \ 4.3;\ \ 12.045$.\par
\pa{b} Change to decimal fractions: $\dfrac25=\dfrac4{10}=0.4$;\par $\dfrac74=\dfrac{175}{100}=1.75$;\par $3\dfrac9{20}=3\dfrac{45}{100}=3.45$;\par $\dfrac{11}8=\dfrac{1\,375}{1\,000}=1.375$.
}
\gap
\begin{probbox}{2}{\relax}
Write as decimal fractions: $\ 0.27;\ \ 0.048;\ \ 3.6;\ \ 15.0091$.
\end{probbox}
\probsol{
\[0.27=\frac{27}{100};\qquad0.048=\frac{48}{1\,000};\qquad3.6=\frac{36}{10};\qquad15.0091=\frac{150\,091}{10\,000}.\]
}
\gap
\begin{probbox}{3}{\relax}
Write the decimal made of:
\begin{cols2}
\pa{a} $6$ ones, $3$ tenths, $4$ hundredths, $1$ thousandth; & \pa{b} $42$ ones, $5$ tenths, $9$ thousandths;\\
\pa{c} $3$ hundreds, $7$ ones, $8$ hundredths; & \pa{d} $0$ ones and $6$ hundredths.
\end{cols2}
\end{probbox}
\probsol{
\pa{a} $6.341$. \qquad \pa{b} $42.509$. \qquad \pa{c} $307.08$. \qquad \pa{d} $0.06$.
}
\gap
\begin{probbox}{4}{\relax}
Order from smallest to largest: $\ 5.34;\ \ 5.043;\ \ 5.4;\ \ 5.304;\ \ 4.99$.
\end{probbox}
\probsol{
$4.99$ has the smallest whole part.\par The other four have whole part $5$; compare tenths: $0<3=3<4$; for the two with $3$ tenths, compare hundredths: $0<4$.
\dapso{$4.99<5.043<5.304<5.34<5.4$.}}
\gap
\begin{probbox}{5}{\relax}
Find the suitable digits:
\begin{cols3}
\pa{a} $7.6\nd{x}2<7.612$ & \pa{b} $3.4\nd{y}5>3.475$ & \pa{c} $5.726<5.7\nd{a}7<5.755$
\end{cols3}
\end{probbox}
\probsol{
\pa{a} Both start $7.6$; compare hundredths: we need $x<1$ (if $x=1$ we would need $2<2$).\par So $x=0$.\par
\pa{b} We need $y>7$ ($y=7$ makes them equal).\par So $y\in\{8;9\}$.\par
\pa{c} $5.7\nd{a}7>5.726$ when $a\ge2$;\par $5.7\nd{a}7<5.755$ when $a\le4$.\par So $a\in\{2;3;4\}$.
}
\gap
\begin{probbox}{6}{\relax}
Find all natural numbers $n$ such that
\begin{cols2}
\pa{a} $0.7<n<3.2$; & \pa{b} $46.85<n<50.1$.
\end{cols2}
\end{probbox}
\probsol{
\pa{a} $n\in\{1;\ 2;\ 3\}$. \qquad \pa{b} $n\in\{47;\ 48;\ 49;\ 50\}$.
}
\gap
\begin{probbox}{7}{\relax}
Find two consecutive natural numbers $m$ and $n$ such that
\begin{cols2}
\pa{a} $m<27.4<n$; & \pa{b} $m<99.9<n$.
\end{cols2}
\end{probbox}
\probsol{
\pa{a} $m=27$, $n=28$. \qquad \pa{b} $m=99$, $n=100$.
}
\gap
\begin{probbox}{8}{\relax}
Find all decimals $x$ with two decimal places such that $\ 0.345<x<0.4$.
\end{probbox}
\probsol{
Write $0.4=0.40$.\par The numbers with two decimal places greater than $0.345$ and less than $0.40$ are:
\[0.35;\ \ 0.36;\ \ 0.37;\ \ 0.38;\ \ 0.39.\]
}
\gap
\begin{probbox}{9}{\relax}
Using the four digits $1,\ 3,\ 6,\ 8$ once each, write the smallest and the largest decimals with two decimal places.
\end{probbox}
\probsol{
The whole part has two digits.\par Smallest: digits in increasing order, $13.68$.\par Largest: in decreasing order, $86.31$.
\dapso{$13.68$ and $86.31$.}}
\gap
\begin{probbox}{10}{\relax}
Calculate in the most convenient way:
\begin{cols2}
$A=0.01+0.02+0.03+\dots+0.09+0.10$ & $B=1.3+1.5+1.7+\dots+19.9+20.1$
\end{cols2}
\end{probbox}
\probsol{
\pa{A} $10$ terms: $A=(0.01+0.10)\times10:2=0.55$.\par
\pa{B} Difference $0.2$; number of terms: $(20.1-1.3):0.2+1=95$.\par $B=(1.3+20.1)\times95:2=1\,016.5$.
}
\gap
\begin{probbox}{11}{\relax}
Work out mentally:
\begin{cols3}
$3.75\times10$ & $0.068\times100$ & $52.1:10$\\
$7.4:100$ & $0.5\times1\,000$ & $216:1\,000$
\end{cols3}
\end{probbox}
\probsol{
\[37.5;\qquad6.8;\qquad5.21;\qquad0.074;\qquad500;\qquad0.216.\]
}
\gap
\begin{probbox}{12}{\relax}
Fill in the decimals:
\begin{cols3}
\pa{a} $3$ m $45$ cm $=\ \dots$ m & \pa{b} $2$ kg $50$ g $=\ \dots$ kg & \pa{c} $4$ km $25$ m $=\ \dots$ km\\
\pa{d} $905$ g $=\ \dots$ kg & \pa{e} $38$ dm$^2=\ \dots$ m$^2$ & \pa{g} $7$ m$^2$ $5$ dm$^2=\ \dots$ m$^2$
\end{cols3}
\end{probbox}
\probsol{
\pa{a} $3.45$ m. \quad \pa{b} $2.05$ kg. \quad \pa{c} $4.025$ km. \quad \pa{d} $0.905$ kg. \quad \pa{e} $0.38$ m$^2$. \quad \pa{g} $7.05$ m$^2$.
\note{Note: $1$ m$^2=100$ dm$^2$, so $5$ dm$^2=0.05$ m$^2$.}
}
\gap
\begin{probbox}{13}{\relax}
The first rope is $3$ m $15$ cm long. The second is $8$ dm longer than the first. The third is $65$ cm shorter than the second. Find the total length of the three ropes in metres.
\end{probbox}
\probsol{
Convert: $3$ m $15$ cm $=3.15$ m; $8$ dm $=0.8$ m; $65$ cm $=0.65$ m.\par
Second rope: $3.15+0.8=3.95$ (m).\par Third rope: $3.95-0.65=3.3$ (m).\par
Total: $3.15+3.95+3.3=10.4$ (m).
\dapso{$10.4$ m.}}
\gap
\begin{probbox}{14}{\relax}
Moving the decimal point of a number one place to the left makes it $16.74$ smaller. Find the number.
\end{probbox}
\probsol{
If the new number is $1$ part, the old one is $10$ parts; the difference is $9$ parts, which is $16.74$.\par
New number: $16.74:9=1.86$.\par The number: $1.86\times10=18.6$.
\dapso{$18.6$.}}
\gap
\begin{probbox}{15}{\relax}
Moving the decimal point of a number one place to the right makes it $48.24$ larger. Find the number.
\end{probbox}
\probsol{
If the old number is $1$ part, the new one is $10$ parts; the difference, $9$ parts, is $48.24$.\par
The number: $48.24:9=5.36$.\par \emph{Check:} $53.6-5.36=48.24$.
\dapso{$5.36$.}}
\gap
\begin{probbox}{16}{\relax}
Write the decimal made of: \ \pa{a} five ones, four tenths, seven hundredths; \ \pa{b} zero ones, three thousandths; \ \pa{c} twelve ones, six hundredths.
\end{probbox}
\probsol{
\pa{a} $5.47$.\par \pa{b} $0.003$.\par \pa{c} $12.06$ (the tenths digit is $0$).
\dapso{$5.47$; $0.003$; $12.06$.}}
\gap
\begin{probbox}{17}{\relax}
What is the value of the digit $7$ in each number? \ \pa{a} $3.472$; \ \pa{b} $70.25$; \ \pa{c} $0.007$.
\end{probbox}
\probsol{
\pa{a} Hundredths: $\dfrac7{100}$.\par \pa{b} Tens: $70$.\par \pa{c} Thousandths: $\dfrac7{1\,000}$.
\dapso{$\dfrac7{100}$; $70$; $\dfrac7{1\,000}$.}}
\gap
\begin{probbox}{18}{\relax}
Write as decimals: \ $\dfrac3{10}$; \ $\dfrac{47}{100}$; \ $\dfrac9{1\,000}$; \ $1\dfrac34$.
\end{probbox}
\probsol{
$\dfrac3{10}=0.3$;\par $\dfrac{47}{100}=0.47$;\par $\dfrac9{1\,000}=0.009$.\par $1\dfrac34=1\dfrac{75}{100}=1.75$.
\dapso{$0.3$; $0.47$; $0.009$; $1.75$.}}
\gap
\begin{probbox}{19}{\relax}
Write as decimals (change the denominator to $10$, $100$ or $1\,000$): \ $\dfrac35$; \ $\dfrac7{20}$; \ $\dfrac9{25}$; \ $\dfrac38$.
\end{probbox}
\probsol{
$\dfrac35=\dfrac6{10}=0.6$;\par $\dfrac7{20}=\dfrac{35}{100}=0.35$.\par $\dfrac9{25}=\dfrac{36}{100}=0.36$;\par $\dfrac38=\dfrac{375}{1\,000}=0.375$.
\dapso{$0.6$; $0.35$; $0.36$; $0.375$.}}
\gap
\begin{probbox}{20}{\relax}
Fill in $>,\ <,\ =$: \ \pa{a} $4.7\ \square\ 4.70$; \ \pa{b} $0.09\ \square\ 0.1$; \ \pa{c} $12.305\ \square\ 12.35$.
\end{probbox}
\probsol{
\pa{a} A zero added at the end of the decimal part changes nothing: $=$.\par \pa{b} Tenths: $0<1$, so $<$.\par \pa{c} Hundredths: $0<5$, so $<$.
\dapso{$=$; $<$; $<$.}}
\gap
\begin{probbox}{21}{\relax}
Order from largest to smallest: $0.8,\ 0.08,\ 0.88,\ 0.808$.
\end{probbox}
\probsol{
With three decimal places: $0.800,\ 0.080,\ 0.880,\ 0.808$.\par Compare: $880>808>800>80$.
\dapso{$0.88>0.808>0.8>0.08$.}}
\gap
\begin{probbox}{22}{\relax}
Write as decimals: \ \pa{a} $5$ m $7$ cm $=\ldots$ m; \ \pa{b} $450$ g $=\ldots$ kg; \ \pa{c} $3$ quintals $5$ kg $=\ldots$ quintals ($1$ quintal $=100$ kg); \ \pa{d} $2$ km $30$ m $=\ldots$ km.
\end{probbox}
\probsol{
\pa{a} $7$ cm $=\dfrac7{100}$ m: $5.07$ m.\par \pa{b} $450$ g $=\dfrac{450}{1\,000}$ kg: $0.45$ kg.\par \pa{c} $5$ kg $=\dfrac5{100}$ quintal: $3.05$ quintals.\par \pa{d} $30$ m $=\dfrac{30}{1\,000}$ km: $2.03$ km.
\dapso{$5.07$; $0.45$; $3.05$; $2.03$.}}
\gap
\begin{probbox}{23}{\relax}
Write as decimals: \ \pa{a} $3$ m$^2$ $25$ dm$^2=\ldots$ m$^2$; \ \pa{b} $7$ dm$^2$ $5$ cm$^2=\ldots$ dm$^2$; \ \pa{c} $4\,500$ m$^2=\ldots$ ha.
\end{probbox}
\probsol{
Each area unit is $100$ times the next smaller one.\par \pa{a} $3.25$ m$^2$.\par \pa{b} $7.05$ dm$^2$.\par \pa{c} $1$ ha $=10\,000$ m$^2$, so $4\,500$ m$^2=0.45$ ha.
\dapso{$3.25$; $7.05$; $0.45$.}}
\gap
\begin{probbox}{24}{\relax}
Work out mentally: \ $3.75\times10$; \ $0.42\times100$; \ $56.3:10$; \ $7.8:100$.
\end{probbox}
\probsol{
Multiplying by $10$, $100$ moves the point $1$, $2$ places right: $37.5$;\par $42$.\par Dividing by $10$, $100$ moves it $1$, $2$ places left: $5.63$;\par $0.078$.
\dapso{$37.5$; $42$; $5.63$; $0.078$.}}
\gap
\begin{probbox}{25}{\relax}
Calculate: \ \pa{a} $12.5+3.75$; \ \pa{b} $20-7.36$; \ \pa{c} $4.5+2.35-1.8$.
\end{probbox}
\probsol{
\pa{a} $12.50+3.75=16.25$.\par \pa{b} $20.00-7.36=12.64$.\par \pa{c} $6.85-1.8=5.05$.
\dapso{$16.25$; $12.64$; $5.05$.}}
\gap
\begin{probbox}{26}{\relax}
Find $x$: \ \pa{a} $x+2.75=6.1$; \ \pa{b} $x-4.6=3.85$.
\end{probbox}
\probsol{
\pa{a} $x=6.1-2.75=3.35$.\par \pa{b} $x=3.85+4.6=8.45$.
\dapso{$x=3.35$; $x=8.45$.}}
\gap
\begin{probbox}{27}{\relax}
Calculate cleverly: $3.25+4.8+6.75+5.2$.
\end{probbox}
\probsol{
$=(3.25+6.75)+(4.8+5.2)$\par $=10+10=20$.
\dapso{$20$.}}
\gap
\begin{probbox}{28}{\relax}
A shop sells sugar on three days: $12.5$ kg on day one, $3.75$ kg more on day two, and $2.4$ kg less than day two on day three. How much sugar does it sell in all?
\end{probbox}
\probsol{
Day two: $12.5+3.75=16.25$ kg.\par Day three: $16.25-2.4=13.85$ kg.\par In all: $12.5+16.25+13.85=42.6$ kg.
\dapso{$42.6$ kg.}}
\gap
\begin{probbox}{29}{\relax}
Find every decimal $x$ with one decimal place such that $3.4<x<3.9$.
\end{probbox}
\probsol{
The whole part is $3$ and the tenths digit is more than $4$ and less than $9$.\par $x=3.5,\ 3.6,\ 3.7,\ 3.8$.
\dapso{$3.5;\ 3.6;\ 3.7;\ 3.8$.}}
\gap
\begin{probbox}{30}{\relax}
Write the largest and the smallest decimal with three different digits, two decimal places and a non-zero whole part.
\end{probbox}
\probsol{
Largest: the biggest digits, in decreasing order: $9.87$.\par Smallest: whole part $1$, then $0$ and $2$: $1.02$.
\dapso{$9.87$ and $1.02$.}}
\dividerbar{Challenge -- Extension}
\begin{probbox}{31}{\stars{1}}
Moving the decimal point of a number two places to the left makes it $40.887$ smaller. Find the number.
\end{probbox}
\probsol{
Moving two places left divides by $100$: new number $1$ part, old number $100$ parts, difference $99$ parts.\par
New number: $40.887:99=0.413$.\par The number: $0.413\times100=41.3$.
\dapso{$41.3$.}}
\gap
\begin{probbox}{32}{\stars{1}}
Two decimals add up to $151.5$. Moving the decimal point of the first number two places left gives the second number. Find both numbers.
\end{probbox}
\probsol{
The first number is $100$ times the second: second $1$ part, first $100$ parts, total $101$ parts.\par
Second number: $151.5:101=1.5$.\par First number: $1.5\times100=150$.
\dapso{$150$ and $1.5$.}}
\gap
\begin{probbox}{33}{\stars{1}}
A decimal, with its point moved one place to the right, is added to the original number; the result is $75.9$. Find the number.
\end{probbox}
\probsol{
If the number is $1$ part, the shifted number is $10$ parts; the sum is $11$ parts.\par
The number: $75.9:11=6.9$.\par \emph{Check:} $69+6.9=75.9$.
\dapso{$6.9$.}}
\gap
\begin{probbox}{34}{\stars{1}}
Replace $a,\ b$ with suitable digits: $\ \nd{1a.b}+\nd{a.b}=20.8$.
\end{probbox}
\probsol{
Note that $\nd{1a.b}=10+\nd{a.b}$.\par So $10+\nd{a.b}+\nd{a.b}=20.8$, i.e.\par $\nd{a.b}\times2=10.8$ and $\nd{a.b}=5.4$.\par
So $a=5$, $b=4$.\par \emph{Check:} $15.4+5.4=20.8$.
}
\gap
\begin{probbox}{35}{\stars{1}}
Consider the sequence $0.5,\ 1,\ 1.5,\ 2,\ \dots,\ 49.5,\ 50$.
\begin{cols2}
\pa{a} How many terms are there? & \pa{b} What is the $37$th term?\\
\pa{c} Find the sum of all terms. & \pa{d} Find the sum of the first $40$ terms.
\end{cols2}
\end{probbox}
\probsol{
\pa{a} $(50-0.5):0.5+1=100$ terms. \qquad \pa{b} $0.5+36\times0.5=18.5$.\par
\pa{c} $(0.5+50)\times100:2=2\,525$. \qquad \pa{d} The $40$th term is $20$; sum: $(0.5+20)\times40:2=410$.
}
\gap
\begin{probbox}{36}{\stars{1}}
Using $0,\ 4,\ 6,\ 8$ once each, how many decimals can be written whose whole-number part and decimal part both have two digits?
\end{probbox}
\probsol{
The number looks like $\nd{ab.cd}$ with $a\ne0$: $a$ has $3$ choices, $b$ has $3$, $c$ has $2$, $d$ has $1$: $3\times3\times2\times1=18$.
\dapso{$18$ numbers.}}
\gap
\begin{probbox}{37}{\stars{2}}
Find the smallest decimal with all digits different whose digit sum is $10$.
\end{probbox}
\probsol{
To be smallest, the whole part must be $0$, and the decimal digits must be as small as possible, in increasing order.\par
$0$ is already used in the whole part, so take $1,\ 2,\ 3$ in turn; what is left is $10-1-2-3=4$ (unused).\par The number is $0.1234$.
\dapso{$0.1234$.}}
\gap
\begin{probbox}{38}{\stars{2}}
A decimal has two decimal places. Removing its decimal point gives a number $1\,221.66$ greater than the original. Find the decimal.
\end{probbox}
\probsol{
Removing the point of a number with two decimal places multiplies it by $100$: old $1$ part, new $100$ parts, difference $99$ parts.\par
The number: $1\,221.66:99=12.34$.
\dapso{$12.34$.}}
\gap
\begin{probbox}{39}{\stars{1}}
Moving the decimal point of a number two places to the right makes it $207.9$ larger. Find the number.
\end{probbox}
\probsol{
The new number is $100$ times the old: old $1$ part, new $100$ parts, difference $99$ parts.\par The number is $207.9:99=2.1$.
\dapso{$2.1$.}}
\gap
\begin{probbox}{40}{\stars{1}}
Two decimals add up to $22.11$. Moving the point of the smaller one place right gives the larger. Find them.
\end{probbox}
\probsol{
The larger is $10$ times the smaller; $11$ parts in all.\par Smaller: $22.11:11=2.01$.\par Larger: $20.1$.
\dapso{$2.01$ and $20.1$.}}
\gap
\begin{probbox}{41}{\stars{1}}
Two numbers differ by $49.5$. Moving the point of the larger one place left gives the smaller. Find them.
\end{probbox}
\probsol{
The larger is $10$ times the smaller; the difference is $9$ parts.\par Smaller: $49.5:9=5.5$.\par Larger: $55$.
\dapso{$5.5$ and $55$.}}
\gap
\begin{probbox}{42}{\stars{1}}
Consider $0.1,\ 0.3,\ 0.5,\ \dots,\ 9.9$. How many terms are there? Find their sum.
\end{probbox}
\probsol{
Step: $0.2$.\par Number of terms: $(9.9-0.1):0.2+1=50$.\par Sum: $(0.1+9.9)\times50:2=250$.
\dapso{$50$ terms; $250$.}}
\gap
\begin{probbox}{43}{\stars{1}}
How many decimals with two decimal places are greater than $2$ and less than $3$?
\end{probbox}
\probsol{
They are $2.01,\ 2.02,\ \dots,\ 2.99$.\par The decimal part runs from $01$ to $99$: $99$ numbers.
\dapso{$99$ numbers.}}
\gap
\begin{probbox}{44}{\stars{1}}
Find every decimal $x$ with one decimal place such that $5.15<x<5.52$.
\end{probbox}
\probsol{
$5.15<5.2$ and $5.5<5.52$.\par $x=5.2,\ 5.3,\ 5.4,\ 5.5$.
\dapso{$5.2;\ 5.3;\ 5.4;\ 5.5$.}}
\gap
\begin{probbox}{45}{\stars{1}}
Calculate quickly: $1.2+2.3+3.4+4.5+5.6+6.7+7.8+8.9$.
\end{probbox}
\probsol{
This list goes up by $1.1$ and has $8$ terms.\par Sum: $(1.2+8.9)\times8:2=10.1\times4=40.4$.
\dapso{$40.4$.}}
\gap
\begin{probbox}{46}{\stars{1}}
Find digits $a$ and $b$ with $\nd{a.b}+\nd{b.a}=7.7$ and $a-b=3$.
\end{probbox}
\probsol{
The sum is $(a+b)$ ones plus $(a+b)$ tenths, i.e.\par $(a+b)\times1.1=7.7$.\par So $a+b=7$, and with $a-b=3$ we get $a=5$, $b=2$.\par Check: $5.2+2.5=7.7$.
\dapso{$a=5$, $b=2$.}}
\gap
\begin{probbox}{47}{\stars{1}}
Moving the point of a decimal one place right and then subtracting the original number gives $31.5$. Find the number.
\end{probbox}
\probsol{
The new number is $10$ times the old; the difference is $9$ times the old.\par The number is $31.5:9=3.5$.
\dapso{$3.5$.}}
\gap
\begin{probbox}{48}{\stars{1}}
Three decimals add up to $11.1$. The first is $1.5$ more than the second, which is $1.2$ more than the third. Find them.
\end{probbox}
\probsol{
The second is $1.2$ more than the third; the first is $1.5+1.2=2.7$ more than the third.\par Three times the third: $11.1-1.2-2.7=7.2$, so the third is $2.4$.\par Second: $3.6$; first: $5.1$.
\dapso{$5.1$; $3.6$; $2.4$.}}
\gap
\begin{probbox}{49}{\stars{2}}
Using $2,\ 5,\ 7$ once each, write every decimal with one decimal place and add them up.
\end{probbox}
\probsol{
There are $6$: $25.7,\ 27.5,\ 52.7,\ 57.2,\ 72.5,\ 75.2$.\par Each digit is in each place $2$ times; the digits add to $14$.\par Sum: $14\times2\times(10+1+0.1)=28\times11.1=310.8$.
\dapso{$310.8$.}}
\gap
\begin{probbox}{50}{\stars{2}}
Writing the block $142857$ over and over after the point gives $0.142857142857\dots$ What is the $100$th digit after the point?
\end{probbox}
\probsol{
Each block has $6$ digits; $100=6\times16+4$.\par After $16$ blocks, the $100$th digit is the $4$th digit of a block: $8$.
\dapso{$8$.}}
\clearpage
\titlehead
\begingroup\setbox0=\vbox{\section*{Three Basic Fraction Problems}}\endgroup
\chapterhead{4}{Three Basic Fraction Problems}

\shead{The three basic problems}
\begin{knowbox}\begin{itemize}[leftmargin=1.4em,itemsep=4pt]
\item \textbf{Problem 1 -- A fraction of a number.} To find $\dfrac mn$ of $A$, compute $\ A\times\dfrac mn\ $ (or $A:n\times m$). \quad Example: $\dfrac25$ of $60$ is $60\times\dfrac25=24$.
\item \textbf{Problem 2 -- A number from the value of its fraction.} If $\dfrac ab$ of a number is $X$, the number is $\ X:\dfrac ab\ $ (or $X:a\times b$). \quad Example: if $\dfrac25$ of a number is $24$, the number is $24:\dfrac25=60$.
\item \textbf{Problem 3 -- The ratio of two numbers.} The ratio of $a$ to $b$ ($b\ne0$) is $\ a:b\ $ or $\ \dfrac ab$, usually in lowest terms. Both quantities must be in \emph{the same unit}. \quad Example: the ratio of $45$ to $20$ is $\dfrac{45}{20}=\dfrac94$.
\end{itemize}\end{knowbox}
\gap
\shead{A fraction of what is left}
\begin{knowbox}
When a problem says ``$\dfrac ab$ \emph{of what is left}'', first find the remaining part (as a fraction of the whole), then take the fraction of it. For instance, using $\dfrac13$ of a journey leaves $\dfrac23$; going on for $\dfrac35$ \emph{of the rest} covers $\dfrac35\times\dfrac23=\dfrac25$ of the journey.\par\smallskip
A \textbf{bar model} shows clearly which part matches the known number.
\end{knowbox}
\gap
\begin{exbox}{\ExWord\ 1}
A garden has an area of $720$ m$^2$. $\dfrac38$ of it grows vegetables, $\dfrac5{12}$ grows flowers, and the rest is paths. Find the area of the paths.
\guide
\textbf{Way 1.} Vegetables: $720\times\dfrac38=270$ (m$^2$).\par Flowers: $720\times\dfrac5{12}=300$ (m$^2$).\par
Paths: $720-(270+300)=150$ (m$^2$).\par
\textbf{Way 2.} The paths take $1-\dfrac38-\dfrac5{12}=\dfrac{24-9-10}{24}=\dfrac5{24}$ of the garden, i.e.\par $720\times\dfrac5{24}=150$ (m$^2$).
\end{exbox}
\gap
\begin{exbox}{\ExWord\ 2}
In the school chess club, $\dfrac35$ of the members play chess, $\dfrac14$ play Chinese chess, and the remaining $12$ play tic-tac-toe. How many members does the club have?
\guide
The $12$ tic-tac-toe players are $\ 1-\dfrac35-\dfrac14=\dfrac{20-12-5}{20}=\dfrac3{20}\ $ of the members.\par
Number of members: $\ 12:\dfrac3{20}=80$.
\end{exbox}
\gap
\begin{exbox}{\ExWord\ 3}
A cyclist rides for three days. On day one she covers $\dfrac13$ of the route, on day two $\dfrac35$ of the remaining distance, and on day three the last $24$ km. How long is the route?
\guide
\begin{barfigC}
\rowlab{-0.15}{0}{Route:}
\blk{0}{2.25}{0}{0.45}{tmBlue}{Day 1}
\blk{2.25}{4.95}{0}{0.45}{tmGreen}{Day 2}
\blk{4.95}{6.75}{0}{0.45}{tmOrange}{$24$ km}
\brcup{2.25}{6.75}{0.42}{left after day 1}
\end{barfigC}
After day one, $1-\dfrac13=\dfrac23$ of the route is left.\par
Day two covers $\dfrac35$ of that, so day three covers $1-\dfrac35=\dfrac25$ of the remainder, i.e.\par $\dfrac25\times\dfrac23=\dfrac4{15}$ of the route.\par
The route is $\ 24:\dfrac4{15}=90$ (km).
\end{exbox}
\gap
\begin{exbox}{\ExWord\ 4}
Three numbers add up to $170$. The first is $\dfrac12$ of the second, and the second is $\dfrac34$ of the third. Find the three numbers.
\guide
Use a common ``part'': the second number is $2$ parts compared with the first and $3$ parts compared with the third.\par Take it as $6$ parts:
\[\text{first}=3\ \text{parts};\quad\text{second}=6\ \text{parts};\quad\text{third}=8\ \text{parts}.\]
Total: $3+6+8=17$ parts; one part is $170:17=10$.\par
The numbers are $30,\ 60,\ 80$.
\end{exbox}
\gap
\begin{exbox}{\ExWord\ 5}
A number plus $45$ equals $\dfrac85$ of the number. Find the number.
\guide
The number is $\dfrac55$ of itself; adding $45$ makes it $\dfrac85$.\par So $45$ is $\dfrac85-\dfrac55=\dfrac35$ of the number.\par
The number is $\ 45:\dfrac35=75$.
\end{exbox}
\dividerbar{Practice}
\begin{probbox}{1}{\relax}
A farm keeps $240$ birds: chickens, ducks and geese. Chickens are $\dfrac38$ of all the birds and ducks are $\dfrac13$. How many geese are there?
\end{probbox}
\probsol{
Chickens: $240\times\dfrac38=90$.\par Ducks: $240\times\dfrac13=80$.\par Geese: $240-90-80=70$.
\dapso{$70$ geese.}}
\gap
\begin{probbox}{2}{\relax}
A school has $1\,200$ pupils. Grade 5 makes up $\dfrac15$ of them, and girls are $\dfrac7{12}$ of Grade 5. How many boys and how many girls are in Grade 5?
\end{probbox}
\probsol{
Grade 5: $1\,200\times\dfrac15=240$ pupils.\par Girls: $240\times\dfrac7{12}=140$.\par Boys: $240-140=100$.
\dapso{$100$ boys and $140$ girls.}}
\gap
\begin{probbox}{3}{\relax}
A runner covers $30$ km in three hours. In the first hour he runs $\dfrac25$ of the distance, in the second hour $\dfrac12$ of the remaining distance. How far does he run in the third hour?
\end{probbox}
\probsol{
First hour: $30\times\dfrac25=12$ (km); left: $30-12=18$ (km).\par
Second hour: $18\times\dfrac12=9$ (km).\par Third hour: $18-9=9$ (km).
\dapso{$9$ km.}}
\gap
\begin{probbox}{4}{\relax}
Mum buys $12$ kg of rice. First she uses $\dfrac13$ of it plus $1$ kg more. Next she uses $\dfrac12$ of what is left plus $\dfrac12$ kg more. How much rice is left?
\end{probbox}
\probsol{
First time: $12\times\dfrac13+1=5$ (kg); left: $12-5=7$ (kg).\par
Second time: $7\times\dfrac12+\dfrac12=4$ (kg); left: $7-4=3$ (kg).
\dapso{$3$ kg.}}
\gap
\begin{probbox}{5}{\relax}
$\dfrac23$ of the pupils in class 5B is $22$ pupils. How many pupils are in class 5B?
\end{probbox}
\probsol{
Pupils in class 5B: $22:\dfrac23=33$.
\dapso{$33$ pupils.}}
\gap
\begin{probbox}{6}{\relax}
After $18$ litres are taken from a can of oil, $\dfrac58$ of the original amount is left. How much oil was there at first?
\end{probbox}
\probsol{
The $18$ litres are $1-\dfrac58=\dfrac38$ of the original.\par At first: $18:\dfrac38=48$ litres.
\dapso{$48$ litres.}}
\gap
\begin{probbox}{7}{\relax}
A number minus $28$ equals $\dfrac35$ of the number. Find the number.
\end{probbox}
\probsol{
$28$ is $1-\dfrac35=\dfrac25$ of the number, so the number is $28:\dfrac25=70$.
\dapso{$70$.}}
\gap
\begin{probbox}{8}{\relax}
This season a field's rice harvest grew by $\dfrac14$ compared with last season, reaching $15$ tonnes. How much was harvested last season?
\end{probbox}
\probsol{
If last season is $1$ (that is $\dfrac44$), this season is $1+\dfrac14=\dfrac54$.\par
Last season: $15:\dfrac54=12$ tonnes.
\dapso{$12$ tonnes.}}
\gap
\begin{probbox}{9}{\relax}
After selling $\dfrac37$ of its boxes of cakes plus $4$ more boxes, a shop has $28$ boxes left. How many boxes did it have at first?
\end{probbox}
\probsol{
Selling only $\dfrac37$ would leave $28+4=32$ boxes, which is $1-\dfrac37=\dfrac47$ of the boxes.\par
At first: $32:\dfrac47=56$ boxes.
\dapso{$56$ boxes.}}
\gap
\begin{probbox}{10}{\relax}
A racing cyclist covers $\dfrac4{15}$ of a route on day one, $\dfrac25$ of the route on day two, and the last $100$ km on day three. How long is the route?
\end{probbox}
\probsol{
$100$ km is $1-\dfrac4{15}-\dfrac6{15}=\dfrac5{15}=\dfrac13$ of the route.\par Route: $100:\dfrac13=300$ km.
\dapso{$300$ km.}}
\gap
\begin{probbox}{11}{\relax}
Two barns hold equal amounts of rice. Taking $\dfrac25$ from the first barn and $\dfrac3{10}$ from the second leaves the first with $12$ tonnes less than the second. How much did each barn hold at first?
\end{probbox}
\probsol{
The first barn keeps $\dfrac35=\dfrac6{10}$ and the second keeps $\dfrac7{10}$ of its original amount.\par
$12$ tonnes is $\dfrac7{10}-\dfrac6{10}=\dfrac1{10}$, so each barn held $12:\dfrac1{10}=120$ tonnes.
\dapso{$120$ tonnes each.}}
\gap
\begin{probbox}{12}{\relax}
Nam has some marbles. Blue ones are $\dfrac13$ of all, yellow ones are $\dfrac34$ of the blue ones, and the other $30$ are red. How many marbles does Nam have?
\end{probbox}
\probsol{
Yellow: $\dfrac34\times\dfrac13=\dfrac14$ of all.\par Red: $1-\dfrac13-\dfrac14=\dfrac5{12}$ of all.\par
All marbles: $30:\dfrac5{12}=72$.
\dapso{$72$ marbles.}}
\gap
\begin{probbox}{13}{\relax}
The first customer buys $\dfrac16$ of the potatoes, the second buys $\dfrac35$ of the rest, and the seller has $18$ kg left. How much did each customer buy?
\end{probbox}
\probsol{
After the first customer $\dfrac56$ is left.\par Finally $\dfrac25$ of that remains: $\dfrac25\times\dfrac56=\dfrac13$ of the potatoes, which is $18$ kg.\par
At first: $18:\dfrac13=54$ kg.\par First customer: $54:6=9$ kg; second: $(54-9)\times\dfrac35=27$ kg.
\dapso{$9$ kg and $27$ kg.}}
\gap
\begin{probbox}{14}{\relax}
Find the ratio of:
\begin{cols3}
\pa{a} $45$ to $20$; & \pa{b} $1.2$ to $3.6$; & \pa{c} $\dfrac23$ m to $75$ cm.
\end{cols3}
\end{probbox}
\probsol{
\pa{a} $\dfrac{45}{20}=\dfrac94$. \qquad \pa{b} $\dfrac{1.2}{3.6}=\dfrac{12}{36}=\dfrac13$.\par
\pa{c} Same unit first: $\dfrac23$ m $=\dfrac{200}3$ cm.\par Ratio: $\dfrac{200}3:75=\dfrac{200}{225}=\dfrac89$.
}
\gap
\begin{probbox}{15}{\relax}
Find $\dfrac34$ of $56$ kg.
\end{probbox}
\probsol{
$56\times\dfrac34=42$ kg.
\dapso{$42$ kg.}}
\gap
\begin{probbox}{16}{\relax}
How many minutes are $\dfrac25$ of an hour?
\end{probbox}
\probsol{
$1$ hour $=60$ minutes.\par $60\times\dfrac25=24$ minutes.
\dapso{$24$ minutes.}}
\gap
\begin{probbox}{17}{\relax}
Class 5C has $40$ pupils, and $\dfrac38$ of them earn top marks. How many is that?
\end{probbox}
\probsol{
$40\times\dfrac38=15$ pupils.
\dapso{$15$ pupils.}}
\gap
\begin{probbox}{18}{\relax}
A roll of cloth is $24$ m long. The shop sells $\dfrac58$ of it. How many metres are left?
\end{probbox}
\probsol{
Sold: $24\times\dfrac58=15$ m.\par Left: $24-15=9$ m.
\dapso{$9$ m.}}
\gap
\begin{probbox}{19}{\relax}
Mr Tu's trip to town is $72$ km. He has covered $\dfrac49$ of it. How far is left?
\end{probbox}
\probsol{
Remaining: $1-\dfrac49=\dfrac59$ of the trip.\par Distance left: $72\times\dfrac59=40$ km.
\dapso{$40$ km.}}
\gap
\begin{probbox}{20}{\relax}
$\dfrac37$ of Hung's money is $42\,000$ dong. How much does he have?
\end{probbox}
\probsol{
$42\,000:\dfrac37=42\,000\times\dfrac73=98\,000$ dong.
\dapso{$98\,000$ dong.}}
\gap
\begin{probbox}{21}{\relax}
Find a number given that $\dfrac56$ of it is $35$.
\end{probbox}
\probsol{
$35:\dfrac56=42$.
\dapso{$42$.}}
\gap
\begin{probbox}{22}{\relax}
An has read $\dfrac25$ of a book and has $90$ pages still to read. How many pages does the book have?
\end{probbox}
\probsol{
$90$ pages are $1-\dfrac25=\dfrac35$ of the book.\par Pages: $90:\dfrac35=150$.
\dapso{$150$ pages.}}
\gap
\begin{probbox}{23}{\relax}
Find the ratio of: \ \pa{a} $15$ to $25$; \ \pa{b} $2$ hours to $40$ minutes; \ \pa{c} $1.2$ m to $30$ cm.
\end{probbox}
\probsol{
\pa{a} $\dfrac{15}{25}=\dfrac35$.\par \pa{b} Same unit: $2$ hours $=120$ minutes; $\dfrac{120}{40}=3$.\par \pa{c} $1.2$ m $=120$ cm; $\dfrac{120}{30}=4$.
\dapso{$\dfrac35$; $3$; $4$.}}
\gap
\begin{probbox}{24}{\relax}
Class 5D has $32$ pupils, $14$ of them girls. Find the ratio of girls to boys, and of girls to the whole class.
\end{probbox}
\probsol{
Boys: $32-14=18$.\par Girls to boys: $\dfrac{14}{18}=\dfrac79$.\par Girls to class: $\dfrac{14}{32}=\dfrac7{16}$.
\dapso{$\dfrac79$; $\dfrac7{16}$.}}
\gap
\begin{probbox}{25}{\relax}
A shop has $360$ kg of rice. On day one it sells $\dfrac14$ of it; on day two, $\dfrac25$ of what is left. How much rice remains?
\end{probbox}
\probsol{
Day one: $360\times\dfrac14=90$ kg; $270$ kg left.\par Day two: $270\times\dfrac25=108$ kg.\par Remaining: $270-108=162$ kg.
\dapso{$162$ kg.}}
\gap
\begin{probbox}{26}{\relax}
Mai has $450\,000$ dong. She spends $\dfrac29$ of it on books and $\dfrac13$ of it on notebooks. How much is left?
\end{probbox}
\probsol{
Books: $450\,000\times\dfrac29=100\,000$ dong.\par Notebooks: $450\,000\times\dfrac13=150\,000$ dong.\par Left: $450\,000-100\,000-150\,000=200\,000$ dong.
\dapso{$200\,000$ dong.}}
\gap
\begin{probbox}{27}{\relax}
After $\dfrac38$ of the water in a tank is used, $750$ litres remain. How much water was there at first?
\end{probbox}
\probsol{
$750$ litres are $1-\dfrac38=\dfrac58$ of the water.\par At first: $750:\dfrac58=1\,200$ litres.
\dapso{$1\,200$ litres.}}
\gap
\begin{probbox}{28}{\relax}
A number plus $\dfrac14$ of itself is $60$. Find the number.
\end{probbox}
\probsol{
The number plus a quarter of it is $\dfrac54$ of the number.\par The number: $60:\dfrac54=48$.
\dapso{$48$.}}
\gap
\begin{probbox}{29}{\relax}
A number minus $\dfrac27$ of itself is $35$. Find the number.
\end{probbox}
\probsol{
What is left is $1-\dfrac27=\dfrac57$ of the number.\par The number: $35:\dfrac57=49$.
\dapso{$49$.}}
\gap
\begin{probbox}{30}{\relax}
An, Binh and Cuong have $84$ marbles. An has $\dfrac13$ of them and Binh $\dfrac37$ of them. How many does each have?
\end{probbox}
\probsol{
An: $84\times\dfrac13=28$.\par Binh: $84\times\dfrac37=36$.\par Cuong: $84-28-36=20$.
\dapso{An $28$, Binh $36$, Cuong $20$ marbles.}}
\dividerbar{Challenge -- Extension}
\begin{probbox}{31}{\stars{1}}
Classes 5A and 5B have $68$ pupils in total, and $\dfrac34$ of class 5A equals $\dfrac23$ of class 5B. How many pupils are in each class?
\end{probbox}
\probsol{
Make the numerators equal: $\dfrac34=\dfrac68$, $\dfrac23=\dfrac69$.\par So $\dfrac18$ of 5A equals $\dfrac19$ of 5B.\par
5A is $8$ equal parts and 5B is $9$; $17$ parts are $68$ pupils, so a part is $4$.\par
5A: $32$ pupils; 5B: $36$ pupils. \ (\emph{Check:} $\dfrac34\times32=24=\dfrac23\times36$.)
}
\gap
\begin{probbox}{32}{\stars{1}}
In a music group, the number of boys is $\dfrac23$ of the number of girls. After $3$ more boys join, the boys are $\dfrac56$ of the girls. How many boys and girls were there at first?
\end{probbox}
\probsol{
The number of girls does not change, so compare everything with it.\par The $3$ new boys are $\dfrac56-\dfrac23=\dfrac16$ of the girls.\par
Girls: $3:\dfrac16=18$.\par Boys at first: $18\times\dfrac23=12$.
\dapso{$12$ boys and $18$ girls.}}
\gap
\begin{probbox}{33}{\stars{1}}
Mr An and Mrs Ha go to the market with $3\,000\,000$ VND in total. Mr An spends $\dfrac59$ of his money and Mrs Ha spends $\dfrac34$ of hers; then they have equal amounts left. How much did each bring?
\end{probbox}
\probsol{
Mr An keeps $\dfrac49$ and Mrs Ha keeps $\dfrac14$ of their money, and these are equal.\par With equal numerators: $\dfrac49$ and $\dfrac14=\dfrac4{16}$.\par
So Mr An's money is $9$ equal parts and Mrs Ha's is $16$; the $25$ parts are $3\,000\,000$ VND, so a part is $120\,000$ VND.\par
Mr An: $1\,080\,000$ VND; Mrs Ha: $1\,920\,000$ VND. (\emph{Check:} each keeps $480\,000$ VND.)
}
\gap
\begin{probbox}{34}{\stars{2}}
Three teams plant trees. Team One plants $\dfrac13$ of all the trees, Team Two plants $\dfrac25$ of the rest, and Team Three plants $30$ trees more than Team Two. How many trees are planted in all?
\end{probbox}
\probsol{
After Team One, $\dfrac23$ of the trees are left.\par Team Two plants $\dfrac25\times\dfrac23=\dfrac4{15}$; Team Three plants $\dfrac23-\dfrac4{15}=\dfrac6{15}$ of all.\par
$30$ trees are $\dfrac6{15}-\dfrac4{15}=\dfrac2{15}$ of all.\par All trees: $30:\dfrac2{15}=225$.
\dapso{$225$ trees.}}
\gap
\begin{probbox}{35}{\stars{2}}
A shop sells all its oranges in $4$ days. Day one sells $\dfrac13$ of what the next three days sell together. Day two sells $\dfrac23$ of day one. Day three sells half of day four. Day four sells $15$ kg more than day one. How many kilograms of oranges were there at first?
\end{probbox}
\probsol{
If day one is $1$ part, the other three days are $3$ parts, $4$ in all: day one sells $\dfrac14$ of the oranges.\par
Day two: $\dfrac23\times\dfrac14=\dfrac16$.\par Days three and four: $1-\dfrac14-\dfrac16=\dfrac7{12}$, split $1:2$, so day four sells $\dfrac7{12}\times\dfrac23=\dfrac7{18}$.\par
$15$ kg is $\dfrac7{18}-\dfrac14=\dfrac{14-9}{36}=\dfrac5{36}$ of the oranges.\par At first: $15:\dfrac5{36}=108$ kg.
\note{Check: the four days sell $27,\ 18,\ 21,\ 42$ kg.}
\dapso{$108$ kg.}}
\gap
\begin{probbox}{36}{\stars{1}}
Mum shares sweets: An gets $\dfrac13$ of them, Binh gets $\dfrac14$, and the remaining $15$ go to little Hoa. How many sweets were there?
\end{probbox}
\probsol{
An and Binh get $\dfrac13+\dfrac14=\dfrac7{12}$ of the sweets.\par So $15$ is $1-\dfrac7{12}=\dfrac5{12}$ of them.\par Sweets: $15:\dfrac5{12}=36$.
\dapso{$36$ sweets.}}
\gap
\begin{probbox}{37}{\stars{1}}
A seller sells $\dfrac25$ of her oranges, then $\dfrac13$ of the rest, then the last $24$. How many oranges did she have?
\end{probbox}
\probsol{
After the first sale $\dfrac35$ are left.\par The second sale is $\dfrac13\times\dfrac35=\dfrac15$ of all the oranges.\par So $24$ is $\dfrac35-\dfrac15=\dfrac25$ of them: $24:\dfrac25=60$.
\dapso{$60$ oranges.}}
\gap
\begin{probbox}{38}{\stars{1}}
Two classes have $76$ pupils. $\dfrac35$ of class A equals $\dfrac23$ of class B. How many pupils are in each class?
\end{probbox}
\probsol{
Make the numerators equal: $\dfrac35=\dfrac6{10}$, $\dfrac23=\dfrac69$.\par $\dfrac6{10}$ of A equals $\dfrac69$ of B, so A is $10$ parts and B is $9$ parts.\par One part: $76:19=4$; A: $40$; B: $36$.
\dapso{$40$ and $36$ pupils.}}
\gap
\begin{probbox}{39}{\stars{1}}
A walker covers $\dfrac14$ of a route on day one, $\dfrac23$ of the rest on day two and the last $15$ km on day three. How long is the route?
\end{probbox}
\probsol{
After day one, $\dfrac34$ is left.\par Day three is $\dfrac13\times\dfrac34=\dfrac14$ of the route.\par Route: $15:\dfrac14=60$ km.
\dapso{$60$ km.}}
\gap
\begin{probbox}{40}{\stars{1}}
Three numbers add up to $140$. The first is $\dfrac23$ of the second, and the second is $\dfrac45$ of the third. Find them.
\end{probbox}
\probsol{
Let the third be $15$ parts; the second is $15\times\dfrac45=12$ parts and the first $12\times\dfrac23=8$ parts.\par $35$ parts in all; one part: $140:35=4$.\par The numbers: $32,\ 48,\ 60$.
\dapso{$32;\ 48;\ 60$.}}
\gap
\begin{probbox}{41}{\stars{1}}
A drum holds oil. After $\dfrac13$ of the oil is taken out, the drum weighs $32$ kg. After another $\dfrac14$ of the original oil is taken out, it weighs $23$ kg. How heavy is the empty drum?
\end{probbox}
\probsol{
$\dfrac14$ of the oil weighs $32-23=9$ kg, so the oil was $36$ kg.\par After the first step $36\times\dfrac23=24$ kg remain.\par Drum: $32-24=8$ kg.
\dapso{$8$ kg.}}
\gap
\begin{probbox}{42}{\stars{1}}
In a class $\dfrac34$ of the pupils like Maths and $\dfrac23$ like Vietnamese; $10$ like both, and everyone likes at least one. How many pupils are there?
\end{probbox}
\probsol{
Adding the two groups counts the $10$ twice: $\dfrac34+\dfrac23=\dfrac{17}{12}$ of the class.\par The extra $\dfrac{17}{12}-1=\dfrac5{12}$ of the class is $10$ pupils.\par Class: $10:\dfrac5{12}=24$.
\dapso{$24$ pupils.}}
\gap
\begin{probbox}{43}{\stars{1}}
A number is increased by $\dfrac15$ of itself; the new number is then decreased by $\dfrac16$ of itself. How does the result compare with the original?
\end{probbox}
\probsol{
After the increase the number is $\dfrac65$ of the original.\par After the decrease it is $\dfrac56$ of that, i.e.\par $\dfrac56\times\dfrac65=1$ times the original.\par The result equals the original.
\dapso{It is equal to the original.}}
\gap
\begin{probbox}{44}{\stars{1}}
Two barns hold $380$ tonnes of paddy. Taking $\dfrac25$ from barn A and $\dfrac13$ from barn B leaves equal amounts. How much did each hold?
\end{probbox}
\probsol{
What is left: $\dfrac35$ of A equals $\dfrac23$ of B, i.e.\par $\dfrac6{10}$ of A equals $\dfrac69$ of B.\par A is $10$ parts, B is $9$; one part: $380:19=20$.\par A: $200$ t; B: $180$ t.
\dapso{$200$ t and $180$ t.}}
\gap
\begin{probbox}{45}{\stars{1}}
A shop sells $\dfrac13$ of a roll of cloth and $2$ m more, then half of the rest and $3$ m more, leaving $7$ m. How long was the roll?
\end{probbox}
\probsol{
Before day two: $(7+3)\times2=20$ m.\par $20+2=22$ m is $\dfrac23$ of the roll.\par Roll: $22:\dfrac23=33$ m.
\dapso{$33$ m.}}
\gap
\begin{probbox}{46}{\stars{1}}
Two people have the same amount of money. One spends $\dfrac35$ of it and the other $\dfrac23$; together they have $330\,000$ dong left. How much did each have?
\end{probbox}
\probsol{
They keep $\dfrac25$ and $\dfrac13$ of their money.\par Together that is $\dfrac25+\dfrac13=\dfrac{11}{15}$ of one person's money, i.e.\par $330\,000$ dong.\par Each had $330\,000:\dfrac{11}{15}=450\,000$ dong.
\dapso{$450\,000$ dong.}}
\gap
\begin{probbox}{47}{\stars{2}}
An gives Binh $\dfrac14$ of his marbles and $3$ more, then gives Cuong half of the rest and $1$ more, keeping $8$. How many did An have?
\end{probbox}
\probsol{
Before giving to Cuong: $(8+1)\times2=18$.\par $18+3=21$ is $\dfrac34$ of the original.\par At first: $21:\dfrac34=28$.
\dapso{$28$ marbles.}}
\gap
\begin{probbox}{48}{\stars{2}}
Lan should find $\dfrac34$ of a number but finds $\dfrac43$ of it instead, getting $14$ too much. Find the number and the correct answer.
\end{probbox}
\probsol{
The wrong result exceeds the right one by $\dfrac43-\dfrac34=\dfrac7{12}$ of the number.\par The number: $14:\dfrac7{12}=24$.\par Correct answer: $24\times\dfrac34=18$.
\dapso{$24$; $18$.}}
\gap
\begin{probbox}{49}{\stars{2}}
The top shelf holds $\dfrac23$ as many books as the bottom shelf. Moving $10$ books up makes the shelves equal. How many books are on each?
\end{probbox}
\probsol{
The total is fixed.\par The top shelf is $\dfrac25$ of it before and $\dfrac12$ after.\par So $10$ books are $\dfrac12-\dfrac25=\dfrac1{10}$ of the total, which is $100$.\par Top: $40$; bottom: $60$.
\dapso{$40$ and $60$ books.}}
\gap
\begin{probbox}{50}{\stars{2}}
An egg seller sells $\dfrac25$ of the eggs in the morning and $\dfrac34$ of the rest in the afternoon, with $30$ left at the end. How many eggs were there?
\end{probbox}
\probsol{
After the morning $\dfrac35$ are left.\par At the end $\dfrac14\times\dfrac35=\dfrac3{20}$ are left, i.e.\par $30$ eggs.\par Eggs: $30:\dfrac3{20}=200$.
\dapso{$200$ eggs.}}
\clearpage
\titlehead
